\documentclass[../main.tex]{subfiles}

\begin{document}
\chapter{Manual de instalación}
\label{C:instalacion}
\lstset{style=asgardstyle, gobble=1}

Este capítulo describe cómo iniciar la aplicación desarrollada ``Super Asgard-V''. La aplicación está disponible en \url{https://jmarquezt.github.io/super-asgard-v/}. Si se desea hacer una instalación local, se proporcionan dos posibilidades: instalación nativa en el servidor o equipo elegidos e instalación utilizando un contenedor Docker. Ambos caminos producen el mismo resultado, la aplicación servida a través de un servidor web y accesible por el navegador.

\section{Requisitos previos}

\begin{itemize}
	\item \textbf{Git}, para clonar el repositorio.
	\item \textbf{Node.js 22} y \textbf{npm} (incluido en Node.js); necesarios para la instalación nativa.
	\item \textbf{Docker Engine} y el plugin \textbf{Docker Compose} (\texttt{docker compose}); necesarios para la instalación con contenedores. No hace falta tener Node.js instalado en el sistema anfitrión en este caso: la imagen de Docker proporcionada ya lo incluye.
	\item Un navegador web moderno (Chrome, Firefox, Edge) con JavaScript habilitado.
\end{itemize}

\section{Instalación nativa en Linux}

Se exponen los pasos necesarios para realizar la instalación en un sistema Linux o en Windows usando el subsistema de Windows para Linux (WSL).

\subsection{Obtener el código fuente}

El Listado~\ref{lst:obtener-codigo} presenta los comandos necesarios para descargar el código fuente de la aplicación en la carpeta deseada.

\begin{lstlisting}[caption={Comandos para obtener el código fuente.},label={lst:obtener-codigo}]
	cd /carpeta-deseada
	git clone https://github.com/jmarquezt/super-asgard-v.git super-asgard-v
	cd super-asgard-v
\end{lstlisting}

\subsection{Instalar dependencias}

Se recomienda utilizar \texttt{npm ci} en lugar de \texttt{npm install}: instala exactamente las versiones fijadas en \texttt{package-lock.json}, garantizando una instalación reproducible (es el mismo comando que usan el despliegue automático con GitHub Actions y el desarrollo). El Listado~\ref{lst:generar-dependencias} presenta el comando necesario para descargar e instalar las dependencias del proyecto.

\begin{lstlisting}[caption={Comando para instalar las dependencias.},label={lst:generar-dependencias}]
	npm ci
\end{lstlisting}

\subsection{Ejecutar en modo desarrollo}

El comando \texttt{npm start} levanta un servidor de desarrollo con recarga automática al modificar el código fuente. El Listado~\ref{lst:servidor-desarrollo} presenta el comando necesario para ejecutar la aplicación en modo desarrollo.

\begin{lstlisting}[caption={Comando para levantar el servidor de desarrollo.},label={lst:servidor-desarrollo}]
	npm start
\end{lstlisting}

La aplicación es accesible desde el navegador en \url{http://localhost:4200}.

\subsection{Generar la build de producción}

\begin{lstlisting}[caption={Comando para compilar la aplicación en modo producción.},label={lst:compilar-produccion}]
	npm run build
\end{lstlisting}

El comando del Listado~\ref{lst:compilar-produccion} genera los ficheros estáticos compilados y optimizados en \path{dist/super-asgard-v/browser/}. Al tratarse de una aplicación estática (sin backend), cualquier servidor web puede publicarlos; por ejemplo, \texttt{http-server} (paquete de npm, se instala una sola vez con \texttt{npm install -g http-server}) o servidores como Nginx o Apache:

\begin{lstlisting}[caption={Comando para iniciar un servidor web.},label={lst:servidor-web}]
	http-server dist/super-asgard-v/browser -p 8080
\end{lstlisting}

Usando \texttt{http-server} como se muestra en el Listado~\ref{lst:servidor-web}, la aplicación es accesible desde el navegador en \url{http://localhost:8080}.

\subsection{Ejecutar la suite de pruebas (opcional)}

El Listado~\ref{lst:ejecutar-pruebas} presenta los comandos que hay que ejecutar para lanzar la suite de pruebas deseada.

\begin{lstlisting}[caption={Comandos para ejecutar la suite de pruebas.},label={lst:ejecutar-pruebas}]
	npm test               # suite completa
	npm run coverage       # suite completa + informe de cobertura
	npm run test:unit      # solo pruebas unitarias
	npm run test:integration  # solo pruebas de integración de componentes
	npm run test:functional   # solo pruebas funcionales (programas ASG completos)
\end{lstlisting}

\section{Instalación con Docker}

El repositorio incluye un \texttt{Dockerfile} multi-etapa y un \texttt{docker-compose.yml} en su raíz, pensados para desplegar la aplicación sin necesidad de tener instaladas las dependencias en la máquina anfitriona.

\subsection{Cómo está construida la imagen}

El \texttt{Dockerfile} define dos etapas:

\begin{enumerate}
	\item \textbf{Compilación}: la imagen (\texttt{node:22-alpine}) copia el código fuente, ejecuta \texttt{npm ci} y compila la aplicación en modo producción con \texttt{ng build --configuration production}, generando \path{dist/super-asgard-v/browser}.
	\item \textbf{Servidor}: la imagen (\texttt{nginx:alpine}) copia únicamente los ficheros estáticos generados en la etapa anterior a \path{/usr/share/nginx/html} y expone el puerto 80. La imagen final no contiene Node.js ni el código fuente, solo un servidor Nginx sirviendo los ficheros compilados.
\end{enumerate}

El fichero \texttt{docker-compose.yml} define un único servicio (\texttt{super-asgard-v}) que arranca esta última imagen y mapea el puerto 80 del contenedor al puerto \textbf{4815} de la máquina anfitriona, con reinicio automático (\texttt{restart: always}).

\subsection{Construir y levantar el contenedor}

Como \texttt{docker-compose.yml} referencia la imagen por nombre (\texttt{super-asgard-v:latest}) sin una sección \texttt{build:}, hay que construirla explícitamente con \texttt{docker build} antes de levantar el servicio la primera vez (y cada vez que el código fuente cambie) tal y como se muestra en el Listado~\ref{lst:compilar-imagen}.

\begin{lstlisting}[caption={Comandos para construir la imagen de la aplicación.},label={lst:compilar-imagen}]
	# 1. Construir la imagen a partir del Dockerfile del repositorio
	docker build -t super-asgard-v:latest .
	
	# 2. Levantar el contenedor en segundo plano
	docker compose up -d
\end{lstlisting}

La aplicación queda disponible en \url{http://localhost:4815}.

\subsection{Otras operaciones habituales}

El Listado~\ref{lst:comandos-docker} presenta una serie de comandos habituales para trabajar con Docker.

\begin{lstlisting}[caption={Comandos para realizar operaciones habituales con Docker.},label={lst:comandos-docker}]
	docker compose logs -f      # ver los logs de Nginx en tiempo real
	docker compose down         # parar y eliminar el contenedor
	docker compose up -d --force-recreate   # reiniciar tras reconstruir la imagen
\end{lstlisting}

\subsection{Cambiar el puerto publicado}

Si el puerto 4815 ya está en uso en la máquina anfitriona, basta con editar el mapeo de puertos en \texttt{docker-compose.yml} (el número de la izquierda es el del anfitrión; el de la derecha, 80, es interno del contenedor y no debe cambiarse) como se muestra en el Listado~\ref{lst:codigo-docker}.

\begin{lstlisting}[caption={Código modificado con otro puerto.},label={lst:codigo-docker}]
	ports:
	- "8080:80"   # la aplicación pasaría a estar en http://localhost:8080
\end{lstlisting}

\subsection{Solución de problemas}

\begin{itemize}
	\item \textbf{\texttt{npm ci} falla con un error de versión de Node}: comprobar con \texttt{node -v} que la versión instalada es 22.x. Herramientas como \texttt{nvm} permiten tener varias versiones de Node conviviendo en el mismo sistema.
	\item \textbf{El puerto 4200, 4815 u 8080 ya está en uso}: liberar el puerto (\texttt{lsof -i :4200} en Linux para ver qué proceso lo ocupa) o publicar la aplicación en otro puerto siguiendo las instrucciones de la sección anterior.
	\item \textbf{Docker no encuentra la imagen \texttt{super-asgard-v:latest}}: falta ejecutar \texttt{docker build} antes de \texttt{docker compose up}: \texttt{docker-compose.yml} no construye la imagen automáticamente, solo la referencia por nombre.
	\item \textbf{Cambios en el código no se reflejan tras reconstruir la imagen}: Docker puede reutilizar capas en caché; para forzar una reconstrucción completa se puede utilizar \texttt{docker build --no-cache -t super-asgard-v:latest .}
\end{itemize}
	
\end{document}